\subsection{根图式}

定义 2. 根图式（或简称图式，如无混淆之虞）是一个三元组 \(D = (M, M_{0}, R)\)，其中：

\((RD_{0})\) M 是一个有限型的自由 Z-模，且子模 \(M_{0}\) 是 M 的直因子；

\((\mathrm{RD}_{\mathrm{I}})\) R 是 M 的一个有限子集；\(\mathrm{R} \cup \mathrm{M}_0\) 生成 Q-向量空间 \(\mathbf{Q} \otimes \mathrm{M}\)；

\((\mathrm{RD}_{\mathrm{II}})\) 对所有 \(\alpha \in \mathrm{R}\)，存在元素 \(\alpha^{\vee}\)（\(\mathrm{M}^{*} = \mathrm{Hom}_{\mathbf{Z}}(\mathrm{M}, \mathbf{Z})\) 中的），使得 \(\alpha^{\vee}(\mathrm{M}_{0}) = 0\)、\(\alpha^{\vee}(\alpha) = 2\)，且 M 的自同态 \(x \mapsto x - \alpha^{\vee}(x)\alpha\) 保持 R 不变。

由第六章 §1 第 1 节，对所有 \(\alpha \in \mathbb{R}\)，元素 \(\alpha^{\vee}\)（\(\mathrm{M}^*\) 中的）由 \(\alpha\) 唯一确定；我们用 \(s_{\alpha}\) 表示 M 的自同态 \(x \mapsto x - \alpha^{\vee}(x)\alpha\)。此外（同前处），\(\mathbf{Q}\)-向量空间 \(\mathbf{Q} \otimes \mathrm{M}\) 是 \(\mathbf{Q} \otimes \mathrm{M}_0\) 与由 R 生成的向量子空间 V(R) 的直和，且 R 是 V(R) 中的一个根系（同前处，定义 1）。

R 的元素称为根图式 D 的根，而元素 \(\alpha^{\vee}\)（\(M^{*}\) 中的）称为逆根。由 M 的诸自同构 \(s_{\alpha}\) 生成的群称为 D 的 Weyl 群，记作 \(\mathrm{W(D)}\)；\(\mathrm{W(D)}\) 的元素在 \(M_{0}\) 上诱导出恒等变换，而在 \(\mathrm{V(R)}\) 上诱导出根系 R 的 Weyl 群中的变换。

例. 1) 对每个有限型的自由 \(\mathbf{Z}\)-模 \(M\)，三元组 \((M, M, \varnothing)\) 是一个根图式。

\begin{enumerate}
\def\labelenumi{\arabic{enumi})}
\setcounter{enumi}{1}
\tightlist
\item
  若 \(D = (M, M_{0}, R)\) 是一个根图式，设 \(M_{0}^{*}\) 是 \(V(R)\) 在 \(M^{*}\) 中的正交补，并设 \(R^{\vee}\) 是 D 的逆根全体所成的集合。则 \(D^{\vee} = (M^{*}, M_{0}^{*}, R^{\vee})\) 是一个根图式，称为 D 的逆图式。对所有 \(\alpha \in R\)，对称 \(s_{\alpha^{\vee}}\)（\(M^{*}\) 的）是对称 \(s_{\alpha}\)（M 的）的反步自同构；映射 \(w \mapsto {}^{t} w^{-1}\) 是从 \(W(D)\) 到 \(W(D^{\vee})\) 的同构。此外，\(V(R^{\vee})\) 可自然地等同于 Q-向量空间 \(V(R)\) 的对偶，此时 \(R^{\vee}\) 等同于 R 的逆根系。
\end{enumerate}

若把 \(\mathbf{M}^*\) 的对偶等同于 M，则 \(\mathrm{D}^{\vee}\) 的逆图式等同于 D。

\begin{enumerate}
\def\labelenumi{\arabic{enumi})}
\setcounter{enumi}{2}
\item
  设 \((\mathfrak{g},\mathfrak{h})\) 是一个分裂约化 \(\mathbf{Q}\)-李代数，\(\mathrm{M} \subset \mathfrak{h}\) 是一个许可格（第八章，§12，第 6 节，定义 1）。设 \(\mathrm{M}_0\) 是 M 中与 \((\mathfrak{g},\mathfrak{h})\) 的根正交的子群，并设 \(\mathrm{R}^\vee\) 是诸 \(H_{\alpha}\)（\(\alpha \in \mathrm{R}(\mathfrak{g},\mathfrak{h})\)）所成的集合。则 \((\mathrm{M},\mathrm{M}_0,\mathrm{R}^\vee)\) 是一个根图式，而 \((\mathrm{M}^*,\mathrm{M}_0^*,\mathrm{R}(\mathfrak{g},\mathfrak{h}))\) 是其逆图式。
\item
  设 V 是 Q 上的一个向量空间，R 是 V 中的一个根系；用 P(R) 表示 R 的权群，用 Q(R) 表示 R 的根权格（第六章，§1，第 9 节）。则 (Q(R), 0, R) 与 (P(R), 0, R) 都是根图式。根图式 (M, M₀, S) 同构于形如 (Q(R), 0, R)（相应地 (P(R), 0, R)）的图式，当且仅当 M 由 S 生成（相应地 M* 由 S \(^{v}\) 生成）。
\end{enumerate}

对 P(R) 的每个包含 Q(R) 的子群 X，(X, 0, R) 是一个根图式；而且，在同构意义下，每个满足 M₀ = 0（换言之，S 生成 M 的一个有限指数子群）的图式 (M, M₀, S) 都由这种方式得到。

根图式 \((\mathrm{M},\mathrm{M}_0,\mathrm{R})\) 称为既约的，如果根系 \(\mathbf{R}\) 是既约的（换言之（第六章，§1，第 4 节），如果关系式 \(\alpha ,\beta \in \mathbb{R}\)、\(\lambda \in \mathbf{Z}\)、\(\beta = \lambda \alpha\) 蕴含 \(\lambda = 1\) 或 \(\lambda = -1\)）。例 1) 与 3) 中的图式都是既约的。

